\begin{corollary}[Proved under stated hypotheses: mixed fixed boundary on surfaces]
\label{D:mixed-fixed-boundary}
Let $S$ be a normal proper connected surface over an algebraically
closed field, and let $Y\subsetneq S$ be closed.
Suppose that every singular point $s\in Y$ at which $S$ is not
locally $\mathbf Q$-factorial is isolated in $Y$ and satisfies
at least one of the following conditions:
\begin{enumerate}
\item it admits a regular resolution with one smooth elliptic
curve as reduced exceptional divisor;
\item $S$ is log canonical at $s$.
\end{enumerate}
Then the fixed-boundary quotient $Q(S,Y)$ of D.1 is topologically
finitely generated.  Singularities outside $Y$ are unrestricted.
In positive characteristic, the omitted locally
$\mathbf Q$-factorial points include all points resolved by
simple normal crossings rational trees, without a log canonical
assumption.
\end{corollary}
\begin{proof}
The characteristic-zero case is already
Proposition~\ref{D:fixed-boundary-positive}, so suppose that
the characteristic is positive.  We retain all the locally
$\mathbf Q$-factorial singular points belonging to $Y$.
We resolve all singularities outside $Y$, and at each remaining
singular point of $Y$ use its single elliptic or cyclic resolution.
The cyclic alternative follows from
Lemma~\ref{D:lc-exceptional-configurations}; its rational-tree
alternative is excluded here by
Proposition~\ref{D:rational-tree-Qfactorial}.
Lemma~\ref{D:Qfactorial-partial-resolution} yields a normal
projective surface $q:X\to S$ with these properties.
Put $U=S\setminus Y$ and $U'=q^{-1}(U)$.

Every connected finite \'etale cover of $U$ remains connected
after pullback to $U'$, by the normal birational argument in
Theorem~\ref{D:surface-mixed-boundary}.
Consider one classified by $Q(S,Y)$.
At a strict transform of a codimension-one component of $Y$,
the associated discrete valuation of the function field is
unchanged, so its pullback is tame there by hypothesis.
Every other divisorial component of $X\setminus U'$ lies above
one of the isolated non-$\mathbf Q$-factorial points of $Y$.
Because that point is isolated in $Y$, the cover is \'etale
on its punctured local neighborhood.  The elliptic theorem
\ref{D:single-curve-residue-tameness}, respectively the lc cusp
theorem~\ref{D:lc-cusp-tame}, therefore makes it tame on all of
these new components as well.
There are no additional boundary divisors: the other exceptional
curves lie over $U$ and belong to $U'$.

Thus the composite homomorphism
$\pi_1(U')\to\pi_1(U)\to Q(S,Y)$ is surjective and kills wild
inertia at every codimension-one component of $X\setminus U'$.
It factors through a surjection
\[
             Q(X,X\setminus U')\longrightarrow Q(S,Y).
\]
The source is finitely generated by the projective case of
Proposition~\ref{D:fixed-boundary-positive}, which proves the result.
\end{proof}
